\section{The Church's Deliberate Vocabulary}

Trent teaches that at the Last Supper Christ identified himself as the priest
foretold according to the order of Melchizedek, offered his Body and Blood
under bread and wine, and commanded the apostles and their successors to offer
them. The sacrifice of the Mass is the same sacrifice as Calvary because the
victim is the same and Christ, now offering through the ministry of priests,
is the same offerer; only the manner of offering differs: bloody on the Cross,
unbloody in the Mass (session XXII, chs.~1--2).

The Second Vatican Council calls the Eucharist a memorial of Christ's death and
resurrection by which the sacrifice of the Cross is perpetuated
(\emph{Sacrosanctum Concilium} 47). The Catechism gathers the tradition:
anamnesis is not bare recollection; Christ's Passover is made present; the
sacrifice remains always actual; the Eucharist represents---that is, makes
present---the Cross and applies its fruit; and Cross and Eucharist are one
sacrifice (CCC 1362--1372).

John Paul II accordingly writes that the Eucharist does not merely evoke the
Passion but makes it sacramentally present, while insisting that Christ's death
and resurrection remain unique events of history
(\emph{Ecclesia de Eucharistia} 11--12). Benedict XVI describes the institution
as the anticipation and sacramental concentration of the sacrifice that would
be carried out on Calvary, and situates Eucharistic memorial within the whole
Paschal Mystery (\emph{Sacramentum Caritatis} 9--11).

This vocabulary is not timid. ``Makes present,'' ``same sacrifice,'' and
``always actual'' say more than psychological remembrance. Yet the Church does
not identify sacramental presence with literal time travel. Her chosen terms
protect both realism and the unrepeatable historical character of the Pasch.

